\documentclass[11pt]{article}
\usepackage[margin=1in]{geometry}
\usepackage{amsmath,amssymb,amsthm}
\usepackage{xurl}
\usepackage{hyperref}
\sloppy
\let\oldus\_ \renewcommand{\_}{\oldus\allowbreak}
\newtheorem{theorem}{Theorem}
\newtheorem{lemma}[theorem]{Lemma}
\theoremstyle{definition}
\newtheorem{definition}[theorem]{Definition}
\newcommand{\R}{\mathbb{R}}

\title{Conjecture 00000002510: tensors of bounded rank do not form a closed set,\\ and best rank-2 approximation fails by diverging components}
\author{Jackmeson1}
\date{}

\begin{document}
\maketitle

\section*{Statement}
\emph{Definition: Best low-rank approximation of tensors. Conjecture: The existence of best
low-rank approximations is governed by closure defects: unlike matrices, the set of
bounded-rank tensors is not closed, and the approximation fails via limiting divergence.}

This is a known theorem of de Silva and Lim. The arXiv abstract of
\url{https://arxiv.org/abs/math/0607647} says: ``we argue that the naive approach to this
problem is doomed to failure because, unlike matrices, tensors of order 3 or higher can fail
to have best rank-r approximations''. We give a self-contained, machine-checked proof of the
statement in the smallest case, real $2\times2\times2$ tensors and rank $2$, together with the
matrix contrast and the general closedness criterion.

\section*{Reading}
We prove the following four claims. Each one is a conjunct of the Lean theorem
\texttt{C2510.conjecture\_2510}.
\begin{enumerate}
\item[(R1)] \emph{Unlike matrices.} For all $m,n,r$, the set of real $m\times n$ matrices of
  rank $\le r$ is closed. Hence every real matrix has a best rank-$\le r$ approximation.
\item[(R2)] \emph{Governed by closure.} For every $r$: every tensor has a best rank-$\le r$
  approximation if and only if the set of tensors of rank $\le r$ is closed.
\item[(R3)] \emph{Bounded-rank tensors are not closed.} The set of $2\times2\times2$ tensors
  of rank $\le 2$ is not closed. The tensor $W$ below has rank $3$ and lies in the closure.
  Its distance to the rank-$\le2$ set is $0$, but no rank-$\le2$ tensor attains it, so $W$
  has no best rank-$\le 2$ approximation.
\item[(R4)] \emph{Limiting divergence.} Take any sequence of decompositions
  $S_n = x_n\otimes y_n\otimes z_n + x_n'\otimes y_n'\otimes z_n'$ with $S_n\to W$. Then the
  norms of \emph{both} summands tend to $+\infty$.
\end{enumerate}
We read ``the set of bounded-rank tensors is not closed'' as: for some tensor space and some
rank bound $r$, the set of rank-$\le r$ tensors is not closed. We do not claim this for every
$r$; for instance, rank $\le 0$ gives the closed set $\{0\}$. Only the four claims above are
proved.

\section*{Conventions and definitions}
\begin{itemize}
\item Tensors are $T:\{0,1\}^3\to\R$ (Lean: \texttt{Fin 2 -> Fin 2 -> Fin 2 -> R}). The norm
  is Mathlib's default sup norm $\|T\|=\max_{i,j,k}|T_{ijk}|$, with $\mathrm{dist}(S,T)=\|S-T\|$.
  Closedness, closure and the divergence statement do not depend on the norm, since all norms
  are equivalent in finite dimension. ``Best approximation'' and the infimum are taken for
  this sup norm.
\item $e_1=(1,0)$ and $e_2=(0,1)$ (indices $0,1$). The outer product is
  $(a\otimes b\otimes c)_{ijk}=a_ib_jc_k$.
\item $T$ has rank $\le r$ (\texttt{HasRankLE T r}) if $T=\sum_{s<r}a_s\otimes b_s\otimes c_s$
  for some vectors $a_s,b_s,c_s\in\R^2$. Also $\mathrm{rankLE}\,r$ is the set of such $T$, and
  $\mathrm{tensorRank}\,T$ is the least such $r$ (an \texttt{sInf} of natural numbers).
\item Matrix rank is Mathlib's \texttt{Matrix.rank}, the dimension of the column space.
\item $y$ is a best approximation of $x$ from $A$ (\texttt{IsBestApprox A x y}) if $y\in A$
  and $\mathrm{dist}(x,y)\le \mathrm{dist}(x,z)$ for every $z\in A$.
\item $W=e_1\otimes e_1\otimes e_2+e_1\otimes e_2\otimes e_1+e_2\otimes e_1\otimes e_1$, so
  $W_{001}=W_{010}=W_{100}=1$ and all other entries are $0$.
\end{itemize}

\section*{Proofs}
\begin{lemma}[\texttt{bestApprox\_forall\_iff\_isClosed}]
Let $A$ be a nonempty subset of a proper metric space $X$. Then every $x\in X$ has a best
approximation from $A$ if and only if $A$ is closed.
\end{lemma}
\begin{proof}
If $A$ is closed, then Mathlib's \texttt{IsClosed.exists\_infDist\_eq\_dist} gives $y\in A$
with $\mathrm{dist}(x,y)=\inf_{z\in A}\mathrm{dist}(x,z)$. Conversely, let $x\in\overline A$
and let $y$ be a best approximation. For each $\varepsilon>0$ there is $z\in A$ with
$\mathrm{dist}(x,z)<\varepsilon$. Hence $\mathrm{dist}(x,y)<\varepsilon$, so $x=y\in A$.
\end{proof}
(R2) follows, because $\R^{2\times2\times2}$ is proper and $0\in\mathrm{rankLE}\,r$.

\begin{theorem}[(R1), \texttt{matrix\_rankLE\_isClosed}, \texttt{matrix\_bestApprox\_exists}]
$\{M\in\R^{m\times n}:\operatorname{rank}M\le r\}$ is closed.
\end{theorem}
\begin{proof}
The map $M\mapsto (v\mapsto Mv)$ from matrices to continuous linear maps is linear on a
finite-dimensional space, hence continuous. The rank of $M$ equals the rank of this linear
map. Mathlib's \texttt{isOpen\_setOfPred\_nat\_le\_rank} says that
$\{f: r+1\le\operatorname{rank}f\}$ is open. Our set is the complement of its preimage. The
existence of best approximations then follows from the Lemma.
\end{proof}

\begin{lemma}[\texttt{W\_eqns\_inconsistent}, \texttt{W\_not\_rankLE\_two}]
$W$ is not a sum of two rank-one tensors.
\end{lemma}
\begin{proof}
Suppose $W=a\otimes b\otimes c+p\otimes q\otimes r$. This gives eight equations
$a_ib_jc_k+p_iq_jr_k=W_{ijk}$. Put $D=a_0p_1-a_1p_0$. Take $p_1$ times the $(0,j,k)$
equation minus $p_0$ times the $(1,j,k)$ equation (Cramer's rule on the first-mode slices).
This gives $Db_jc_k=p_1W_{0jk}-p_0W_{1jk}$, that is,
$Db_0c_0=-p_0$, $Db_0c_1=Db_1c_0=p_1$ and $Db_1c_1=0$. Since
$(Db_0c_1)(Db_1c_0)-(Db_0c_0)(Db_1c_1)=0$ identically, we get $p_1^2=0$. In the same way,
$a_0$ times the $(1,j,k)$ equation minus $a_1$ times the $(0,j,k)$ equation gives
$Dq_0r_0=a_0$, $Dq_0r_1=Dq_1r_0=-a_1$ and $Dq_1r_1=0$, so $a_1^2=0$. Then the $(1,0,0)$
equation $a_1b_0c_0+p_1q_0r_0=1$ reads $0=1$. Every step is a \texttt{linear\_combination}
in Lean.
\end{proof}
Hence $W$ is not of rank $\le 0,1,2$, using the padding lemma \texttt{HasRankLE.mono}. The
three-term definition of $W$ is a rank-$3$ decomposition (\texttt{W\_hasRankLE\_three}), so
$\mathrm{tensorRank}\,W=3$ (\texttt{tensorRank\_W}).

\begin{theorem}[(R3), \texttt{Tseq\_mem}, \texttt{Tseq\_tendsto}, \texttt{W\_mem\_closure},
\texttt{rankLE\_two\_not\_isClosed}, \texttt{infDist\_W}, \texttt{W\_no\_best\_approx}]
Let $t_n=1/(n+1)$ and $u_n=e_1+t_ne_2$. Then
$T_n=(n+1)\,u_n\otimes u_n\otimes u_n-(n+1)\,e_1\otimes e_1\otimes e_1$ has rank $\le2$ and
$T_n\to W$. Consequently $\mathrm{rankLE}\,2$ is not closed,
$\inf_{S\in\mathrm{rankLE}\,2}\|W-S\|=0$, and $W$ has no best rank-$\le2$ approximation.
\end{theorem}
\begin{proof}
An entry-by-entry check (\texttt{Tseq\_eq}) gives $T_n=W+t_nP+t_n^2Q$, where
$P=e_1\otimes e_2\otimes e_2+e_2\otimes e_1\otimes e_2+e_2\otimes e_2\otimes e_1$ and
$Q=e_2\otimes e_2\otimes e_2$. Since $t_n\to0$, we get $T_n\to W$, so $W$ lies in the closure.
But $W\notin\mathrm{rankLE}\,2$ by the Lemma, so the set is not closed. The infimum is $0$
by \texttt{Metric.mem\_closure\_iff\_infDist\_zero}. A best approximation $S$ would satisfy
$\|W-S\|\le\|W-T_n\|\to0$, so $W=S\in\mathrm{rankLE}\,2$, which is impossible.
\end{proof}

\begin{theorem}[(R4), \texttt{max\_summand\_tendsto}, \texttt{summand\_tendsto}]
Let $x_n,y_n,z_n$ (summand $0$) and $x'_n,y'_n,z'_n$ (summand $1$) be vectors in $\R^2$
with $S_n=x_n\otimes y_n\otimes z_n+x'_n\otimes y'_n\otimes z'_n\to W$. Then
$\|x_n\otimes y_n\otimes z_n\|\to\infty$ and $\|x'_n\otimes y'_n\otimes z'_n\|\to\infty$.
\end{theorem}
\begin{proof}
First, $\max(\|\text{summand }0\|,\|\text{summand }1\|)\to\infty$. Otherwise there are $C$
and a subsequence $\varphi$ on which both summand norms are $<C$. In the sup norm,
$\|a\|\|b\|\|c\|\le\|a\otimes b\otimes c\|$ (\texttt{norm\_mul\_le\_norm\_outer3}). So each
summand can be rewritten with factors $(\|b\|\|c\|)a$, $b/\|b\|$, $c/\|c\|$ (or $0,0,0$ if
$b=0$ or $c=0$), all of norm $\le\max(C,1)$ (\texttt{rebalance}). By Bolzano--Weierstrass
(Mathlib's \texttt{tendsto\_subseq\_of\_bounded}), a further subsequence of these factor
pairs converges. By continuity of $(a,b,c)\mapsto a\otimes b\otimes c$, the limit is a
two-term decomposition of $\lim S_{\varphi(\psi(n))}=W$. This contradicts the Lemma.
Second, the triangle inequality gives
$\|\text{summand }s\|\ge\max(\cdot,\cdot)-\|S_n\|$, and eventually $\|S_n\|<\|W\|+1$.
\end{proof}
The Wikipedia article \url{https://en.wikipedia.org/wiki/Tensor_rank_decomposition} describes
this phenomenon (the summand norms tending to infinity) and says: ``It is sometimes called
the problem of diverging components.''

\section*{Lean}
The file \texttt{lean/Conjecture2510/Basic.lean} (namespace \texttt{C2510}, only
\texttt{import Mathlib}) proves the main theorem \texttt{C2510.conjecture\_2510}. It is the
conjunction of (R1) for all $m,n,r$, (R2) for all $r$, non-closedness of
$\mathrm{rankLE}\,2$, $W\in\overline{\mathrm{rankLE}\,2}$, $\mathrm{tensorRank}\,W=3$,
$\mathrm{infDist}(W,\mathrm{rankLE}\,2)=0$, the absence of a best approximation, and (R4)
for both summands. \texttt{lake build} succeeds, and \texttt{\#print axioms
C2510.conjecture\_2510} reports only \texttt{propext}, \texttt{Classical.choice} and
\texttt{Quot.sound}. There is no \texttt{sorry} and no \texttt{native\_decide}.

\end{document}
